\section{Virtual Embodiment}
\label{sec:background_virtual_embodiment}

Virtual reality can do more than surround the user with a simulated environment. It can also provide the user with a virtual body, body parts, tools, or task-specific apparatus that appear to occupy the space of action. In such cases, the user does not only perceive a scene from the outside, but may experience a relation between their own body and the body or action represented in the virtual environment. This relation is commonly discussed under the concept of virtual embodiment.
% TODO cite: general virtual embodiment / sense of embodiment overview.

Virtual embodiment is usually not treated as a single, indivisible feeling. Instead, it is often described through multiple related components, especially body ownership, sense of agency, and self-location. Body ownership concerns the experience that a body or body part belongs to oneself. Sense of agency concerns the experience of causing or controlling an action. Self-location concerns where the user experiences themselves to be located in relation to a body or space.
% TODO cite: Kilteni et al. sense of embodiment framework.
% TODO cite: source distinguishing ownership, agency, and self-location.

This distinction is important because the components of embodiment can diverge. A user may feel that they control a virtual hand because it moves in response to their own action, while still not strongly feeling that the hand is part of their body. Conversely, a virtual body may appear spatially aligned with the user and therefore support self-location, while delayed or mismatched movement may reduce agency. Embodiment is therefore better understood as a structured experience in which different cues can support different aspects of the relation between the physical and virtual body.

Many embodiment studies build on the logic of body illusion paradigms. In the rubber hand illusion, for example, synchronous visual and tactile stimulation can lead participants to report ownership over an artificial hand. Virtual hand and virtual body studies extend this logic into immersive environments, where the virtual body can respond to the user's movement and may be experienced as part of the user's body representation.
% TODO cite: Botvinick and Cohen rubber hand illusion.
% TODO cite: virtual hand illusion / virtual body ownership studies.
% TODO cite: Slater / Sanchez-Vives / body ownership in VR where appropriate.

For this thesis, the most relevant aspect of virtual embodiment is not the mere presence of a virtual avatar, but the relation between the user's real movement and the movement shown in VR. In indoor rowing, the user performs a real, physically effortful action on an ergometer. If a virtual body or rowing apparatus is shown, the system must decide how this physical action is represented. The visible movement may closely correspond to the user's actual handle and body motion, or it may be transformed into a visually more sport-specific boat and oar action.

This makes VR rowing a useful case for examining embodiment in a tool-mediated and exercise-based context. The user is not only looking at a virtual body, nor only moving a hand in empty space. They are seated on a physical machine, holding a real handle, pushing against foot stretchers, and performing a rhythmic full-body movement. The virtual representation must therefore mediate between the felt physical apparatus and the visual action shown in the virtual environment.

The central concern of this thesis is how this mediation affects embodiment-related experience. An ergometer-congruent representation may support embodiment because the seen action can remain close to the user's felt movement. A boat-centric representation may support the plausibility and meaning of the sport scenario, but may require a greater transformation between physical action and visual movement. The following sections therefore examine the components of embodiment and the role of visuomotor congruence in more detail.
